\documentclass[border=8pt]{standalone}
\usepackage[utf8]{inputenc}
\usepackage{nacal}
\usepackage{tikz}
\usepackage{tikz-3dplot}
\usetikzlibrary{calc,angles,positioning,intersections,quotes,decorations.markings,arrows.meta,decorations.pathreplacing}
\usepackage{tkz-euclide}
\def\Datos{green!40!gray}
\def\Ajuste{blue!50!red}
\def\Residuo{black!25!gray}
\def\Regresor{blue!75!white}
\def\Base{orange!80!black}
\def\Sistem{blue!50!black}
\tikzset{vector/.style={-{Stealth[length=3mm, width=2mm]},very thick},
         basevec/.style={-{Stealth[length=2.2mm, width=1.6mm]},thick,\Base},
         guide/.style={dashed,thin,gray}}

\begin{document}
  \tdplotsetmaincoords{62}{12}
  \begin{tikzpicture}[scale=2.4, tdplot_main_coords]
    \coordinate (O)  at (0,0,0);
    \coordinate (V)  at (1.0,0.9,0.85);   % v = y - ybar
    \coordinate (P)  at (1.0,0.9,0);      % proyección sobre E
    \coordinate (Px) at (1.0,0,0);
    \coordinate (Py) at (0,0.9,0);
    % el subespacio E
    \fill[gray!10,draw=gray!50,very thin]
      (-0.4,-0.35,0) -- (2.2,-0.35,0) -- (2.2,1.7,0) -- (-0.4,1.7,0) -- cycle;
    \node[gray!70!black,anchor=north east,font=\small,inner sep=2pt] at (2.2,1.7,0) {$\mathcal E$ \ ($\dim\mathcal E=k-1$)};
    % el subespacio R (vertical)
    \draw[gray!60,thin] (O) -- (0,0,1.35) node[anchor=south,font=\small,gray!70!black] {$\mathcal R$ \ ($\dim\mathcal R=n-k$)};
    % base adaptada
    \draw[basevec] (O) -- (0.5,0,0) node[anchor=north,font=\small] {$\boldsymbol b_1$};
    \draw[basevec] (O) -- (0,0.5,0) node[anchor=west,font=\small] {$\boldsymbol b_2$};
    \draw[basevec] (O) -- (0,0,0.5) node[anchor=east,font=\small] {$\boldsymbol b_k$};
    % coordenadas en E
    \draw[guide] (P) -- (Px);
    \draw[guide] (P) -- (Py);
    \node[font=\scriptsize,gray!70!black,anchor=north west,inner sep=1pt] at (Px) {$c_1$};
    \node[font=\scriptsize,gray!70!black,anchor=north east,inner sep=1pt] at (Py) {$c_2$};
    % triángulo
    \fill[\Ajuste!12,opacity=0.5] (O) -- (P) -- (V) -- cycle;
    \tkzMarkRightAngle[size=.10,gray,very thin,fill=lightgray!70,opacity=0.6](O,P,V)
    \draw[vector,\Ajuste] (O) -- (P)
      node[pos=0.55,anchor=north west,font=\small,inner sep=1pt]
      {$\boldsymbol{\mathop{\widehat y}}-\boldsymbol{\mathop{\overline y}}$};
    \draw[vector,\Residuo] (P) -- (V)
      node[pos=0.5,anchor=west,font=\small,inner sep=2pt]
      {$\boldsymbol{\mathop{\widehat e}}$};
    \draw[vector,\Datos] (O) -- (V)
      node[pos=0.5,sloped,above,font=\small,inner sep=1.5pt]
      {$\boldsymbol v=\boldsymbol y-\boldsymbol{\mathop{\overline y}}$};
    \node[font=\small,anchor=south west] at (-0.4,1.35,0.9) {$\mathcal L(\boldsymbol 1)^\perp=\mathcal E\oplus\mathcal R$};
  \end{tikzpicture}
\end{document}
